\section{Orbital Velocities}

With the initial separation $r = 125$ and star mass $M = 1000$ established, we can characterise all relevant velocity thresholds for the planet's orbit. These thresholds determine the qualitative shape of the trajectory and inform the simulation's random velocity selection.

\subsection{Circular Orbital Velocity}

The circular orbital velocity is the speed at which the gravitational centripetal acceleration exactly provides the centripetal force required to maintain a circular path:

\begin{equation}
v_c = \sqrt{\frac{GM}{r}} = \sqrt{\frac{1 \times 1000}{125}} = \sqrt{8} = 2\sqrt{2} \approx 2.83
\end{equation}

At this speed, the planet would trace a perfect circle of radius $r = 125$ around the star.

\subsection{Escape Velocity}

The escape velocity is the minimum speed needed for the planet to escape the star's gravitational field entirely, reaching infinity with zero remaining kinetic energy:

\begin{equation}
v_e = \sqrt{\frac{2GM}{r}} = \sqrt{\frac{2 \times 1000}{125}} = \sqrt{16} = 4
\end{equation}

As expected, $v_e = \sqrt{2}\,v_c \approx 1.414 \times 2.83 = 4$, confirming the general relationship between circular and escape velocities.

\subsection{Velocity Classification Table}

Based on these two critical speeds, all possible orbital behaviours can be classified as follows:

\begin{center}
\begin{tabular}{l c}
\toprule
Speed $v$ & Orbital Behaviour \\
\midrule
$v = 0$ & Planet falls directly into the star \\
$0 < v < 2\sqrt{2}$ & Elliptical orbit (sub-circular) \\
$v = 2\sqrt{2} \approx 2.83$ & Circular orbit \\
$2\sqrt{2} < v < 4$ & Elliptical orbit (super-circular) \\
$v = 4$ & Parabolic trajectory (marginally unbound) \\
$v > 4$ & Hyperbolic escape \\
\bottomrule
\end{tabular}
\end{center}

\vspace{0.5cm}

The two elliptical regimes differ in their eccentricity. Below the circular velocity the orbit is elongated inward (periapsis is closer to the star than the starting position); above it the orbit is elongated outward (apoapsis is further than the starting position).

\subsection{Velocity Selection in the Simulation}

Rather than fixing a single initial speed, the simulation uses a uniform random draw from the interval $[0.5,\; v_e]$. This is implemented in \texttt{physics::choose\_vel} using the Mersenne Twister engine for reproducibility and quality randomness:

\begin{equation}
v_0 \sim \mathcal{U}(0.5,\; 4)
\end{equation}

The lower bound of $0.5$ prevents near-zero velocities that would cause an immediate near-collision. Any value below $v_e = 4$ guarantees a bound orbit; values above $v_e$ are excluded by the distribution's upper bound, so the planet cannot escape. In practice, each run of the simulation produces a qualitatively different orbit — from a highly elliptical path to one that approaches circular — making it straightforward to observe the full range of Keplerian behaviour by simply re-running the program.

For the worked numerical examples in the subsequent sections, a representative speed of $v_0 = 2.8$ is used, placing the orbit in the super-circular elliptical regime just below the escape threshold.